\chapter{Proposed numerical formulation and solution procedure}
\label{ch:numerics}

This chapter specifies the discretization and solution procedure for the model of
Chapters~\ref{ch:equations} and~\ref{ch:interface} at the level required for an independent
implementation review. No code is presented and none is yet written (Appendix~B); the specification
is intended to be realisable with standard finite-element components in a general-purpose
environment.

\section{Spatial discretization}
The fluid velocity and pressure are approximated by Taylor--Hood elements (continuous piecewise
quadratic velocity, continuous piecewise linear pressure), satisfying the discrete inf--sup
condition for the Stokes limit of \eqref{eq:momentum} \cite{hood1974navier,brezzi1991mixed}; both
temperature fields and the solid displacement use the same quadratic family. The fluid energy
equation \eqref{eq:fluidenergy} is strongly advection-dominated: with
$\mathrm{Pe}=\mathrm{Re}_{D_h}\mathrm{Pr}=348$ the cell P\'eclet number exceeds the Galerkin
oscillation threshold on any mesh affordable at the required near-wall resolution, so the
temperature transport is stabilized by a consistent streamline-upwind (SUPG) term with the
standard residual-based coefficient; oscillation-free behaviour of the computed interface flux is
part of the acceptance criteria of test V-B (Chapter~\ref{ch:validation}). The meshes of the
fluid and solid subdomains are constructed \emph{conforming} at the interface: nodes and edges of
the two meshes coincide on $\Ga_{\mathrm{if}}$. This choice is deliberate: interface quantities are
then transferred by identical nodal values, the interface quadrature is exact to the geometric
approximation order, and no mortar or projection operator -- each of which would introduce its own
consistency error in the conjugate flux condition -- is required. Isoparametric quadratic boundary
edges are used on the deformable boundary and are updated with the same map as the mesh
displacement, so that the geometric approximation of $\Om_f$ stays at the approximation order of
the field variables.

\section{Weak formulation and monolithic interface cancellation}
\label{sec:coupling}
Multiplying by test functions and integrating by parts, the coupled steady problem reads: find
$\bu\in V_g$, $p\in Q$, $T_f\in S_f$, $\bd\in V_s$, $T_s\in S_s$ such that
\begin{align}
&\hat a_f(\bu,\mathbf v)+b(\mathbf v,p)+c(\bu;\bu,\mathbf v)
=\int_{\Ga}\sigma_f(\bu,p)\bn\cdot\mathbf v\,\mathrm ds
+\int_{\Ga_{\mathrm{out}}}\sigma_f\bn\cdot\mathbf v\,\mathrm ds +\langle\text{data}\rangle,
\qquad b(\bu,q)=0\ \ \forall\,q, \label{eq:weakf}\\
&\rho c_p\,a_T(T_f,\chi)+a_{kf}(T_f,\chi)
=-\int_{\Ga}k_f\nabla T_f\cdot\bn\,\chi\,\mathrm ds, \label{eq:weaktf}\\
&\int_{\Om_s^0}\bsig(\bd,\theta):\beps(\boldsymbol\varphi)\,\mathrm dX
=\int_{\Ga}\bsig\bn\cdot\boldsymbol\varphi\,\mathrm ds, \label{eq:weaks}\\
&a_{ks}(T_s,\chi)=\int_{\Ga_h^0}q^{''}\chi\,\mathrm ds
-\int_{\Ga}k_s\nabla T_s\cdot\bn\,\chi\,\mathrm ds , \label{eq:weakts}
\end{align}
for all test functions, with
$\hat a_f(\bu,\mathbf v)=2\mu\int_{\Om_f}\mathbf D(\bu):\mathbf D(\mathbf v)\,\mathrm dx$,
$b(\mathbf v,p)=-\int_{\Om_f}p\nabla\cdot\mathbf v\,\mathrm dx$, and
$\sigma_f=-p\bI+2\mu\mathbf D(\bu)$ the fluid Cauchy stress. The interface mechanism is worth
stating exactly, because it is the load-bearing sign argument of the coupling. After integration
by parts, the fluid residual contains the interface term
$-\int_{\Ga}\sigma_f\bn\cdot\mathbf v\,\mathrm ds$ (with $\bn$ the outward normal of $\Om_f$),
while the solid residual contains $+\int_{\Ga}\sigma_s\bn\cdot\boldsymbol\varphi\,\mathrm ds$
(the solid's outward normal on $\Ga$ is $-\bn$). Under the traction equality
\eqref{eq:iftraction}, $\sigma_s\bn=\sigma_f\bn$, and with identical test functions on a
conforming interface mesh the two terms enter the assembled residual with opposite signs and
cancel identically; likewise the flux terms $-\int_{\Ga}k_f\nabla T_f\cdot\bn\,\chi\,\mathrm ds$
and $+\int_{\Ga}k_s\nabla T_s\cdot\bn\,\chi\,\mathrm ds$ cancel under \eqref{eq:ifflux}. The
monolithic statement therefore contains \emph{no} modelled interface quantity, which removes one
class of coupling error by construction. Three conditions are required for the cancellation to
survive discretization, and all three are part of the specification: (i) the single-normal
convention of Section~\ref{sec:signs}; (ii) a conforming interface mesh, so that traction and
temperature transfer without projection; and (iii) identical quadrature and the same
isoparametric representation of the curved P2 interface on both sides, so that the two boundary
integrals are evaluated on the same discrete surface --- this quadrature-consistency audit is
part of tests V-B and V-D. The outflow term in \eqref{eq:weakf} is the do-nothing natural
condition of Section~\ref{sec:bcs} and does not cancel.

\section{Solution procedure for the steady problem}
The domain map $\widehat{\mathbf w}$ is defined as the solution of a pseudo-elastic problem on
$\Om_f^0$ with coefficient graded to stiffen small elements near the interface, Dirichlet data
$\widehat{\mathbf w}=\bd$ on $\Ga_{\mathrm{if}}^0$ and $\widehat{\mathbf w}=\mathbf 0$ on rigid
boundaries -- the standard solid-based mesh-motion device of deformable-domain computations
\cite{lesoinne1996geometric}. The steady coupled state is obtained by the following fixed-point
iteration on the geometry:
\begin{enumerate}\itemsep1pt
\item solve the fluid subproblem on the current $\Om_f$ by exact Newton linearization of
\eqref{eq:weakf} (the advection derivative is linearized analytically);
\item transfer pressure-dominated traction and interface temperature requirements to the solid;
\item solve the thermoelastic subproblem \eqref{eq:weaks}--\eqref{eq:weakts} at the reduced
(frozen-force) von K\'arm\'an level specified by Section~\ref{sec:deriv:beamcolumn} (linear
thermoelasticity augmented by the geometric-stiffness contribution of the current membrane
state; the membrane force is re-evaluated from the computed state after each solve, so the
consistency update \eqref{eq:nupdate} is recovered at the fixed point; one direct solve per
iteration);
\item solve the mesh-extension problem, update $\Om_f$, and recompute all geometry-dependent
data (normals, Jacobians, quadrature);
\item under-relax the interface displacement with a documented factor $\omega$
(tuned once and reported per run);
\item iterate to a smallness of the interface displacement increment,
$\|\bd^{(k+1)}-\bd^{(k)}\|_\infty\le10^{-6}\,w_{\max}$
(achieved values reported per run; the earlier figure $10^{-8}H_0$ sat below achievable
fixed-point accuracy and is superseded), together with residual smallness of all four
subproblem statements.
\end{enumerate}
The escalation trigger to the monolithic path is numeric: escalation occurs when
$N_\theta^c/N_{\mathrm{cr}}>0.5$ at any iterate, or when the geometry increment stagnates
(four consecutive iterations with reduction of the increment by less than a factor of two). At
the baseline ($N_\theta^c/N_{\mathrm{cr}}=0.275$) the fixed point remains the specified path.
If it nevertheless stagnates or shows loss of convergence near the buckling margin,
the recorded escalation is a monolithic Newton solve including consistent shape-derivative terms;
pseudo-transient continuation is the fallback. Linear systems are factorized directly (the
two-dimensional problem size at the target meshes is of order $10^4$--$10^5$ degrees of freedom);
iterative solvers are a scaling option, not a baseline assumption.

\section{Mesh convergence and quality controls}
\label{sec:quality}
Three nested uniformly refined meshes are used per configuration; observed orders of convergence
and grid-convergence indices are reported for the monitored quantities mean Nusselt number, maximum
interface deflection, and pressure drop. A configuration enters the reported set only if the
fine-to-medium difference is within a pre-registered tolerance ($1\%$, tightened to $0.5\%$ for any
headline figure). In addition to the residuals, every computed run must satisfy, and tabulate, the
following conservation and consistency checks, which follow directly from the formulation:
(i) \emph{global enthalpy balance},
$|\int_{\Ga_h}q^{''}\mathrm ds-\dot m'c_p(T_{b,\mathrm{out}}-T_{b,\mathrm{in}})|\le10^{-3}Q'$,
evaluated on the deformed geometry;
(ii) \emph{interface flux-jump norm}, $\|[k\partial_nT]\|_{L^2(\Ga)}/\|q^{''}\|\le10^{-3}$;
(iii) \emph{load-case consistency}: mechanism attribution is performed by separate solutions
(thermal only, traction only, both) rather than by adding fields; exact additivity of the linear
model is verified on a detuned case with $N_\theta^c/N_{\mathrm{cr}}\le0.05$
(Section~\ref{sec:coupling}, Chapter~\ref{ch:validation});
(iv) \emph{mean-stress identity} (\ref{eq:meanstress}a,b) pointwise within interpolation accuracy;
(v) \emph{regime guards} A-1--A-6 of Chapter~\ref{ch:dimensionless} recomputed from the computed
fields.

\section{Post-processing definitions}
The bulk temperature $T_b(x)$ is mass-flux-weighted; the wall heat flux is
$q_w(x)=-k_f\nabla T_f\cdot\bn|_\Ga$; the local heat-transfer coefficient and Nusselt number are
\begin{equation}
h(x)=\frac{q_w(x)}{T_w(x)-T_b(x)},\qquad \mathrm{Nu}(x)=\frac{2H_0h(x)}{k_f},
\end{equation}
with $\mathrm{Nu}$ singular where $T_w\to T_b$ near the start of the heated segment; that singular
behaviour is documented and plotted rather than suppressed. The areal mean is taken over the heated
length. Pressure drop is measured between area-averaged pressures on the two rigid end sections.
Solid post-processing exports complete Cauchy components (never the mechanical part alone) and
forms $\sigma_m$ from (\ref{eq:meanstress}a) with the identity (\ref{eq:meanstress}b) exercised as
a runtime assertion.

\section{Computational feasibility note}
All prescribed components (Taylor--Hood assembly, conforming interface mesh, direct sparse
factorization, mesh-extension mapping, fixed-point loop) are standard and require no specialised
CFD infrastructure; the dominant implementation risk is not algorithmic but metrological: whether
the checks (i)--(v) are exercised, reported, and used as exclusions, as required in
Chapter~\ref{ch:validation}. The verification of the baseline formulation against an independently
implemented model (an open-source finite-element platform) is part of the validation methodology
(Section~\ref{sec:replication}).
